\documentclass[tikz,border=12pt]{standalone}
\usepackage{fontspec}
\usepackage[mathrm=sym]{unicode-math}
\setmainfont{texgyrepagella-regular.otf}[
    BoldFont = texgyrepagella-bold.otf,
    ItalicFont = texgyrepagella-italic.otf,
    BoldItalicFont = texgyrepagella-bolditalic.otf
]
\setmathfont{texgyrepagella-math.otf}[AutoFakeBold = 2.5]
\usepackage{tikz}
\usetikzlibrary{arrows.meta}

\begin{document}
\begin{tikzpicture}[>=Stealth, line cap=round, line join=round]

  % Axes
  \draw[thick, ->] (0, 0) -- (7.2, 0) node[below=2pt, font=\Large] {$s$};
  \draw[thick, ->] (0, 0) -- (0, 5.2) node[left=2pt, font=\Large] {$V(s)$};

  % Potential curve V(s)
  % Ground state at (0, 2.8), peak at (2.0, 4.2), drop to bottom near (6.0, 0.1)
  \draw[very thick, line width=1.4pt, smooth]
    plot coordinates {
      (0, 2.8) (0.6, 2.88) (1.2, 3.4) (1.6, 3.95) (2.0, 4.2) (2.4, 4.05)
      (2.8, 3.4) (3.3, 2.2) (3.8, 1.2) (4.4, 0.5) (5.2, 0.18) (6.4, 0.08)
    };

  % Dashed horizontal lines for reference
  \draw[thick, dashed] (2.0, 4.2) -- (6.3, 4.2);
  \draw[thick, dashed] (0, 2.8) -- (6.3, 2.8);

  % Barrier height E_B
  \draw[{Latex}-{Latex}, thick] (1.8, 2.8) -- (1.8, 4.2);
  \node[right=2pt, font=\Large] at (1.8, 3.5) {$E_B$};

  % Barrier magnitude: ~ 5-10 MeV
  \draw[{Latex}-{Latex}, thick] (5.6, 2.8) -- (5.6, 4.2);
  \node[right=3pt, font=\normalsize] at (5.6, 3.5) {$\sim 5\text{--}10\ \mathrm{MeV}$};

  % Energy release: several hundred MeV
  \draw[{Latex}-{Latex}, thick] (5.6, 0.1) -- (5.6, 2.8);
  \node[right=3pt, align=left, font=\normalsize] at (5.6, 1.45) {several\\hundred $\mathrm{MeV}$};

\end{tikzpicture}
\end{document}
